% !TEX root = tars_report.tex

\documentclass[9pt,a4paper,twoside]{utils/tau-class/tau}
\usepackage[english]{babel}
\renewcommand{\familydefault}{\sfdefault} % use the sans-serif (abstract) font for the whole document
\usepackage{parskip} % no first-line indent; separate paragraphs with vertical space instead
\usepackage{graphicx}
\graphicspath{{figures/}}
% Visible placeholder box for figures whose image is not yet available.
% To insert the real picture, uncomment the \includegraphics line in each figure
% and drop the image into a figures/ folder next to this .tex.
\newcommand{\figplaceholder}[1]{%
    \setlength{\fboxsep}{6pt}%
    \fbox{\parbox[c][3.2cm][c]{0.9\linewidth}{\centering\itshape [\,figure placeholder --- #1\,]}}}

%----------------------------------------------------------
% TITLE
%----------------------------------------------------------

\journalname{Synthesis Project}
\title{TARS: A Translator Assignment Recommendation System for iDISC}

%----------------------------------------------------------
% AUTHORS, AFFILIATIONS AND PROFESSOR
%----------------------------------------------------------

% TODO: replace with the full names of every team member.
% TODO: replace the XXXXXXX placeholders with each member's real NIU.
\author{Team 9 --- Joan Alcaraz\textsuperscript{1}, Hermes Barreiro\textsuperscript{2}, Marc Bell-Lloc\textsuperscript{3}, Diana Dragan\textsuperscript{4}, Adriana M\'endez\textsuperscript{5}, and Eric P\'erez\textsuperscript{6}\par\vspace{4pt}
{\normalfont\footnotesize NIU: \textsuperscript{1}1746392 \quad \textsuperscript{2}1708343 \quad \textsuperscript{3}1747203 \quad \textsuperscript{4}1807526 \quad \textsuperscript{5}1808502 \quad \textsuperscript{6}1724064}}

%----------------------------------------------------------
% FOOTER INFORMATION
%----------------------------------------------------------

\institution{Universitat Aut\`onoma de Barcelona}
\footinfo{Synthesis Project}
\theday{}
\leadauthor{Team 9}
\course{Synthesis Project}

%----------------------------------------------------------
% ABSTRACT AND KEYWORDS
%----------------------------------------------------------

\begin{abstract}
This project builds a recommendation system that helps the project managers (PMs) at iDISC assign the right translator to each incoming task. The company runs a 24/7 operation and handles about 70,000 tasks per year, and the decision of who works on each task is made manually, a process that consumes roughly 30,000 PM hours annually and depends on each manager's personal memory and judgment. Our system takes the description of an incoming task --- the client, the language pair, the time window, the deadline and the task type --- and returns a ranked shortlist of translators who can realistically take it on, together with the reasons behind each suggestion.

The recommendation system works in two stages. First, a constraint engine filters out every translator who cannot do the task because of a hard limitation: the wrong language pair, no availability, a rate above the client's budget, a quality history below the client's minimum, or no experience in the relevant field. When no perfect candidate exists, the engine relaxes the requirements gradually and in a controlled order, mirroring how a human PM compromises in practice. Second, a scoring step ranks the remaining candidates along the dimensions that matter --- quality, punctuality, cost and experience --- with the emphasis shifting according to what each client values most. We implemented two rankers for this: a transparent rule-based weighted score and a machine-learning model.

This report describes the problem and its business context, the data we worked with and what we learned from analysing it, the technical design of our solution, the business value it offers, the difficulties we encountered and how we approached them, the way we organised ourselves as a team, and finally an assessment of how the solution helps people, where it falls short, and what we would do differently if we started again.
\end{abstract}

\keywords{recommendation system, constraint satisfaction, translator assignment, localization, ranking, gradient boosting, decision support}

\begin{document}
    \maketitle
    \thispagestyle{firststyle} \tauabstract
    \tableofcontents

%----------------------------------------------------------

\section{Introduction}

\subsection{Context}

Translation companies operate as a marketplace between clients who need content translated and a pool of translators with different skills, languages, rates, schedules and track records. The company behind this project manages a very large volume of work --- around 70,000 tasks per year --- across many language pairs, many clients and many task types, and it does so around the clock. Behind every one of those tasks sits a small but real decision: who should do it?

\subsection{The Problem}

Choosing a translator for a task is a constrained matching problem with several objectives. The translator must:

\begin{itemize}
    \item \textbf{cover the language pair} of the task (for example, English into Spanish);
    \item \textbf{be available} during the time window and have enough capacity to finish before the deadline;
    \item \textbf{fit the client's budget} --- their hourly rate must be acceptable given what the client pays;
    \item \textbf{meet the client's quality requirements} --- some clients impose a minimum quality level;
    \item ideally, \textbf{have experience} in the client's industry or in the specific type of task.
\end{itemize}

On top of these constraints, different clients have different priorities. Some are most sensitive to price, some to quality, and some to meeting the deadline. A good assignment is not simply the ``best'' translator in the abstract, it is the translator who best fits \emph{this} task for \emph{this} client under \emph{these} priorities, without wasting a good translator on a routine job or putting an underqualified person on a hard one.

Today this matching is done manually. For each task, a PM relies on experience to decide who is suitable, and then assigns the work. At the company's scale, this is:

\begin{itemize}
    \item \textbf{Expensive} --- about 30,000 PM hours per year are spent on assignment decisions that are largely repetitive.
    \item \textbf{Inconsistent} --- two PMs may handle the same task differently, and decisions reflect individual habits rather than the company's full history of what has actually worked.
    \item \textbf{Hard to scale} --- as volume grows, the manual approach does not.
\end{itemize}

\subsection{Objective}

Our objective was to build a system that, given a new task, \textbf{recommends the translators who best fit it} while doing so quickly, consistently, transparently, and on the basis of the company's real historical data. The system is explicitly designed to \emph{support} the PM rather than replace them: it produces a ranked, explained shortlist, and the PM keeps the final decision. This keeps a human accountable for the outcome and makes the tool realistic to use in practice.

\section{The Data}

Understanding the data was both the starting point and one of the hardest parts of the project, so we describe it in some detail here.

\subsection{Sources}

The company provided four datasets (originally four sheets of a single workbook):

\begin{itemize}
    \item \textbf{Task history (\texttt{Data})} --- the central log of past tasks. Each row records the project and task identifiers, the responsible PM team, the planned start and delivery times, the task type, the source and target languages, the assigned translator, a series of workflow timestamps (when the task was pre-assigned, when the translator was told to start, when they began working, when they delivered, when the PM received and closed it), the hours worked, the translator's rate, the total cost, a quality evaluation score, and a four-level classification of the client's industry (sector, industry group, industry, sub-industry).
    \item \textbf{Schedules (\texttt{Schedules})} --- each translator's working hours (start and end of their workday) and which days of the week they work.
    \item \textbf{Clients (\texttt{Clients})} --- each client's selling price per hour, their minimum acceptable quality score, and a ``wildcard'' field stating which requirement (quality or price) may be relaxed when no perfect match is available.
    \item \textbf{Translator rates and pairs (\texttt{TranslatorsCost+Pairs})} --- for each translator, which language pairs they cover and their hourly rate for each.
\end{itemize}

\subsection{Scale and Shape of the Data}

Our analysis was carried out on the full historical dataset. Its scale is substantial (Table~\ref{tab:data-scale}). In total the history covers nearly \textbf{one million tasks} spread across about \textbf{13,692 projects}, handled by \textbf{4 PM teams}, performed by \textbf{983 translators}, for \textbf{2,645 clients}, across \textbf{13 task types}.

\begin{table}[H]
    \centering
    \caption{Size of the four source datasets.}
    \label{tab:data-scale}
    \footnotesize
    \begin{tabular}{p{0.50\columnwidth}cc}
        \toprule
        \textbf{Table} & \textbf{Rows} & \textbf{Columns} \\
        \midrule
        Task history & 997,933 & 29 \\
        Schedules & 983 & 11 \\
        Clients & 2,645 & 6 \\
        Translator rates and pairs & 4,688 & 7 \\
        \bottomrule
    \end{tabular}
\end{table}

The data is highly concentrated in several ways, which has direct consequences for the system:

\begin{itemize}
    \item \textbf{Task types} --- proofreading (about 43\%) and translation (about 40\%) together make up roughly 83\% of all work. The remaining types (post-editing, miscellaneous, engineering, language lead, management, DTP, test, spotcheck, training, and a few others) form a long tail.
    \item \textbf{Languages} --- the work is mostly from English: English is the source language in about 97\% of tasks. Then the target, Spanish (Iberian) is about 49\% of tasks, followed by Galician, Catalan, Spanish (Latin America) and Basque.
    \item \textbf{Language pairs} --- the pair English $\rightarrow$ Spanish (Iberian) represents about 47\% of all tasks. The translation work is therefore mostly a few pairs, and a lot of rare ones.
    \item \textbf{Clients} --- the top two clients alone account for about 43\% of all tasks, and the top five for roughly 56\%. Two industry sectors (Information Technology and Communication Services) together represent about 70\% of the work.
    \item \textbf{Translators} --- workload is very unevenly distributed. The most active translator alone did about 6.5\% of all tasks, and the top twenty translators together did close to half of all the work.
\end{itemize}

Every task is scored on quality on a 0--10 scale, with a mean of about 7.0 and a standard deviation of about 1.5. Translator hourly rates range from 8 to 62, with a typical value around 16--17 on the task records. Clients' selling prices range from 20 to 50 (median 35), and the selling price is, on average, roughly double the translator cost rate, which is where the company's margin comes from.

\subsection{Data Quality}

The historical data, being real and accumulated over more than a decade, contained a number of inconsistencies that we had to identify and handle. The most important findings were:

\begin{itemize}
    \item \textbf{Out-of-order timestamps} --- about 3.7\% of tasks (around 36,800 rows) had a start time recorded as \emph{later} than the end time. The other workflow steps (pre-assigned $\rightarrow$ ready $\rightarrow$ working $\rightarrow$ delivered) were consistent.
    \item \textbf{Zero-value tasks} --- about 7\% of tasks (around 69,000 rows) recorded zero hours worked \emph{and} zero cost, mostly in proofreading.
    \item \textbf{Outliers} --- hours worked ranged up to 725 in a single record and cost up to about 14,000, far above the typical values (median hours around 0.3, median cost around 5).
    \item \textbf{Same-language tasks} --- about 0.16\% of tasks had the same source and target language, which is not a normal translation.
    \item \textbf{Impossible dates} --- isolated records dated 1956 and 2050 clearly are not correct.
    \item \textbf{Encoding artefacts} --- some names and language labels carried character-encoding problems (for example, accented characters rendered incorrectly).
    \item \textbf{Missing values} --- these were, by contrast, almost negligible: only a handful of rows were missing cost, rate or timestamps, and there were no exact duplicate rows in any table.
\end{itemize}

A reassuring finding was that the four datasets join together almost perfectly: essentially every task in the history matched a client, a schedule and a rate record, with only five translator/language-pair combinations failing to match (and those were traced to spelling variants rather than genuinely missing data). This meant we could combine the sources with confidence.

\subsection{Relationships in the Data}

Beyond describing the data, we studied how the variables relate to each other, since this directly informs what a recommendation model can and cannot do.

\begin{itemize}
    \item \textbf{Cost does not buy quality.} The correlation between a translator's rate and the quality score of their work is essentially zero (about $-$0.01). Paying more does not lead to better results.
    \item \textbf{Quality is hard to predict.} None of the available features carry much signal about the quality of a future task. The best predictors are a translator's \emph{prior average quality} (on similar task types, language pairs, or overall), but even these are weak.
    \item \textbf{Quality and punctuality are largely independent.} Being on time and producing high-quality work are almost unrelated (correlation about 0.02 at the task level). This is an important finding: it means the system should treat these as two separate dimensions rather than assuming a ``good'' translator is good on both at once.
    \item \textbf{Some engineered features were redundant or degenerate.} Several candidate features turned out to be perfectly correlated with each other (for example, two different ways of computing margin), or constant across all rows, so they provided no information.
\end{itemize}

\subsection{New Translators}

At the beginning, we also worried about new entities with no history. But we found that this was not really the issue: there were essentially \textbf{no translators with zero history}. The real problem was \textbf{sparse history}, a few hundred translators had only a handful of scored tasks, which makes their estimates unreliable. We also found that a translator's quality estimate tends to stabilise after only a small number of tasks (a median of about four), though some take many more. This changed our perspective from how to handle the absence of a translator's record, to how much to trust a track record with few tasks.

\subsection{The Five Recommendation Dimensions}

From all of the above, we settled on five dimensions that capture what makes a translator a good fit for a task, and which the system is built around:

\begin{enumerate}
    \item \textbf{Experience} --- accumulated hours working for the client, the client's sector or industry, and the task type.
    \item \textbf{Quality} --- the translator's historical quality evaluation scores.
    \item \textbf{Punctuality} --- how reliably the translator delivers on time.
    \item \textbf{Cost} --- the translator's rate against the client's budget (the margin).
    \item \textbf{Availability} --- whether the translator's schedule fits the task's time window and deadline.
\end{enumerate}

\section{Technical Design}

\subsection{Overview}

The system is a pipeline that turns the raw company data into a ranked, explained recommendation for a given task. It separates offline data preparation (done once) from the online recommendation flow (run for each task), so the live response is fast. The main components are a data-preparation stage, a feature-building stage, a constraint engine, a ranking step, an output formatter, a backend service, a front-end interface, and an evaluation script.

\subsection{System Architecture}

TARS is organised in three layers with clear boundaries between them, shown in Figure~\ref{fig:architecture}.

\textbf{The recommendation core} is a set of Python modules (built on \texttt{pandas}) that contain all of the logic described in this section: data loading, cleaning and feature building, the constraint engine, the two rankers, and the output formatter. It is written to be independent of how it is called, so exactly the same code runs from a command-line script and from the web service.

\textbf{The backend API} is a \texttt{FastAPI} service that sits between the core and the front-end. It exposes a \texttt{/recommendations} endpoint that takes a task's parameters, runs the full pipeline, and returns the ranked shortlist together with the requirements the system derived for the task; alongside it, a set of read-only endpoints (\texttt{/translators}, \texttt{/clients}, \texttt{/tasks}, \texttt{/schedules}) let the interface browse the underlying data with pagination, served from a small SQLite database built from the source tables. The front-end only ever talks to these endpoints, never to the recommendation code directly.

\textbf{The web front-end} is a single-page application that project managers use in the browser; it calls the API over HTTP and renders the recommendations and the supporting data.

\textbf{Orchestration.} For development the front-end and the backend are containerised and started together with a single \texttt{docker compose} command, and a bootstrap script prepares the data and the Python environment, so the whole system can be brought up from a clean checkout.

\begin{figure*}[t]
    \centering
    \figplaceholder{system architecture diagram (data core $\rightarrow$ FastAPI backend $\rightarrow$ web front-end) --- replace with \texttt{figures/architecture.png}}
    % \includegraphics[width=\linewidth]{architecture.png}
    \caption{High-level architecture of TARS: an offline data pipeline produces the processed tables that the recommendation core reads; a FastAPI backend exposes the core to a React front-end used by the project managers.}
    \label{fig:architecture}
\end{figure*}

\subsection{Data Preparation and Feature Building}

The raw workbook is first converted into clean per-table CSV files, and then into a set of processed tables read at run time. Preparation cleans the task history (parsing timestamps, computing how long each task took, deriving an on-time flag, and dropping inconsistent rows) and builds a per-translator statistics table. For each translator this brings together, from all four sources: their rate for each language pair; their average quality; their on-time score; their experience by task type; and their experience by industry at three levels --- sub-industry, industry, and industry group. This hierarchical view of experience is what lets the engine reason about expertise at different levels of specificity. We keep the raw data untouched and write all transformations into separate folders, so the process can be re-run from scratch, and a single bootstrap script sets up the environment and builds these tables for a new contributor.

\subsection{The Constraint Engine}

The core of the system decides which translators are eligible for a task. It first looks up the client's requirements (budget, minimum quality, wildcard preference) and derives the relevant industries from the client's history and the delivery window from the dates. It then applies hard constraints that remove translators who cannot do the task: wrong language pair, rate above budget, quality below the client's minimum, on-time score below threshold, no availability or insufficient daily capacity to meet the deadline, and no experience in the required field.

Some task types change these rules, reflecting real business policy: \textbf{test} tasks ignore the budget and select on quality; \textbf{training} tasks drop the quality and experience requirements; a \textbf{translation followed by proofreading} relaxes the quality threshold slightly; and \textbf{proofreading} or \textbf{spotcheck} tasks must go to a different and more experienced translator than the one who did the previous step.

\subsection{Graceful Relaxation and Splitting}

In a marketplace this constrained, it frequently happens that no translator satisfies every hard requirement. Rather than return ``no match'', the engine relaxes the requirements step by step, in a deliberate order, stopping as soon as candidates appear: first it looks for a specialist at the narrowest (sub-industry) level; then it broadens the required expertise to the industry and industry-group levels; then it applies the client's wildcard, relaxing only the single dimension the client cares least about (lowering minimum quality, raising the budget by 25\%, or easing the on-time and scheduling requirements); and finally it relaxes globally as a last resort. This ordering encodes the idea that we would rather find a less specialised translator than violate the client's stated priority. When relaxation still fails --- typically because the task is too large for one person in the time available --- the engine splits the work, either across days for one translator or across several translators working in parallel. Crucially, whenever the system relaxes something, it records what it relaxed, so the recommendation can be explained honestly.

\subsection{Ranking: Rule-Based Scorer}

Once the engine produces a set of eligible translators, they are ranked. Our first ranker is a transparent rule-based score out of 100, combining four normalised components --- quality, reliability, cost margin and experience --- with weights that shift according to the client's wildcard (Table~\ref{tab:weights}).

\begin{table}[H]
    \centering
    \caption{Ranking weights by client priority (wildcard).}
    \label{tab:weights}
    \footnotesize
    \begin{tabular}{p{0.40\columnwidth}cccc}
        \toprule
        \textbf{Mode} & \textbf{Qual.} & \textbf{Rel.} & \textbf{Margin} & \textbf{Exp.} \\
        \midrule
        Balanced (default) & 40 & 30 & 20 & 10 \\
        Quality-sensitive & 60 & 20 & 10 & 10 \\
        Price-sensitive & 20 & 20 & 50 & 10 \\
        Deadline-sensitive & 20 & 60 & 10 & 10 \\
        Test task & 80 & 0 & 0 & 20 \\
        \bottomrule
    \end{tabular}
\end{table}

A notable design choice is what we call \textbf{resource conservation}. Outside explicitly quality-driven cases, the scorer does not simply reward the highest quality and most experience; it rewards a good-enough fit, so a translator who comfortably clears the client's requirements scores higher than one who massively exceeds them. The reasoning is practical: putting the company's most expert translators on routine tasks wastes them when they could be reserved for demanding work. This is a deliberate business stance rather than a pure ``best translator wins'' approach, and it is one of the points we would want to validate carefully with the company.

\subsection{Ranking: Machine-Learning Model}

Our second ranker is a gradient-boosted regression model that predicts an ``efficiency'' score for each eligible candidate from features such as the client's budget and required quality, the translator's quality, on-time record and rate, and the gaps between them. The intention was the same ``good fit, not overkill'' idea expressed as a learned function rather than fixed weights. Both rankers consume the same eligible set, so the constraint logic is shared and only the final ordering differs.

\subsection{Output and Explainability}

Whichever ranker is used, the system returns the top recommendations (by default the top ten) as a table in which each row carries not just the translator and the score, but the underlying values: their rate for the language pair, their daily capacity, their average quality, their on-time score, and the relevant experience figure. This is central to the product: the PM does not just see \emph{who} is recommended, but \emph{why}, and can therefore trust or overrule the suggestion with full information.

\subsection{The Backend Service}

To make the system usable from a web interface, we built a backend service that exposes it over HTTP. It offers an endpoint that takes a task's parameters and returns the ranked recommendations, plus a set of endpoints that let the interface browse the underlying data --- translators, clients, tasks and schedules --- with pagination. The recommendation results are returned in a clean structured form, including the requirements the system derived for the task, so the front-end can display the full rationale.

\subsection{The Web Interface}

The front-end is a single-page web application aimed at project managers, built with React, Vite and TypeScript. It provides a dashboard and dedicated views for assigning a task, comparing translators, and browsing the translators, clients and tasks in the system (Figure~\ref{fig:webui}). A project manager fills in the task parameters, submits them, and gets back the ranked shortlist with the per-dimension values behind each recommendation. The interface communicates with the backend only through the API described above, keeping a clean separation between the user interface and the recommendation logic, so that either side can evolve independently.

\begin{figure}[ht]
    \centering
    \figplaceholder{screenshot of the web interface (task form and ranked recommendations) --- replace with \texttt{figures/web-ui.png}}
    % \includegraphics[width=\linewidth]{web-ui.png}
    \caption{The project-manager web interface, showing a task and its ranked translator recommendations.}
    \label{fig:webui}
\end{figure}

\subsection{Evaluation}

To judge whether the recommendations are actually good, we wrote a backtesting script. It takes historical tasks, hides the translator who was really chosen, runs the full system on the task, and checks where the real translator appears in the recommended list. From this we can compute concrete measures: how often the real translator is in the top $N$ (the hit rate), how often they appear among the eligible candidates at all (recall), and their average rank when they do (Figure~\ref{fig:results}). This gives an objective, data-grounded way to track whether changes to the system make it better or worse, rather than relying on intuition.

\begin{figure}[ht]
    \centering
    \figplaceholder{backtest results (for example, hit rate at top-$N$) --- replace with \texttt{figures/backtest.png}}
    % \includegraphics[width=\linewidth]{backtest.png}
    \caption{Backtesting results: how often the translator actually chosen for a historical task appears among the system's top recommendations.}
    \label{fig:results}
\end{figure}

\subsection{Reproducibility and Deployment}

The whole project is set up to be reproducible. A single bootstrap script creates a clean environment, installs the dependencies, and builds the data tables the system needs, so that a fresh copy of the repository can be made to run with one command. The web interface and the backend can be brought up together for development, and the data flow --- from raw files, to intermediate CSVs, to the final processed tables the system reads --- is documented and scripted end to end.

\section{Business (B2B) Perspective}

Technically, TARS is a recommendation pipeline; commercially, it is a decision-support tool for project managers that targets one of the company's largest hidden costs. Its value falls along four lines. \textbf{Speed and cost}: replacing the manual search with an instant ranked shortlist removes the bulk of the 30,000 PM hours spent on assignment, letting PMs handle more tasks and focus on decisions that genuinely need judgment. \textbf{Consistency and institutional memory}: every task is evaluated the same way and on the basis of the company's entire history, removing variation between managers and capturing knowledge that would otherwise live only in individuals' heads. \textbf{Client-aware flexibility}: the wildcard mechanism means a price-sensitive and a quality-sensitive client asking for the same task get appropriately different recommendations, automatically. \textbf{Trust through explainability}: every recommendation comes with the figures behind it, and the system is honest about when it has relaxed a requirement, which is what makes the difference between a tool PMs use and a black box they ignore. Taken together, this positions TARS not as an automation that removes the PM, but as a force multiplier that keeps a human firmly in control --- which is also what makes it realistic to roll out in a real operation.

\section{Difficulties and Teamwork}

\subsection{Messy Real-World Data}

By far the biggest challenge was the data. A decade of records contained the inconsistencies described above, none of them obvious until we looked, and each forcing a decision about whether to drop, correct, or flag the affected rows. Our approach was to make the data-quality analysis the first piece of work, write it up clearly, and treat its conclusions as the shared foundation everyone else built on, so the rest of the team did not keep rediscovering the same problems.

\subsection{Working at Scale}

Nearly a million rows was enough to cause practical problems: some analysis ran out of memory computing statistics over the whole dataset at once. This pushed us toward a more disciplined design --- compute heavy features once and store them, sample where a full pass was not necessary, and keep the expensive preparation strictly separate from the fast live step.

\subsection{Deciding What to Optimise For}

A more conceptual difficulty was deciding what ``the best translator'' even means. Our analysis showed that quality is hard to predict and unrelated to cost, that punctuality is more predictable and driven by availability, and that the two are essentially independent. These findings pushed us toward treating the dimensions separately and made us appreciate that a clear, explainable rule-based approach was at least as defensible as a more complex model.

\subsection{Teamwork}

We organised the work with an issue tracker, breaking the project into clear areas --- data analysis and feature creation, the constraint engine, the ranking models, the backend, and the front-end --- and assigning owners to each. We used separate branches so people could work in parallel without breaking one another's code, and kept a shared log of key decisions (including how we used AI tools) so everyone worked from the same understanding. The main friction was the ordinary one of integration: making sure independently developed components agreed on the same formats. We addressed this by agreeing on a single documented data flow and a clear contract between the front-end and the backend, so each part could be developed independently against a stable interface.

\section{Conclusion}

\subsection{How Our Solution Helps People}

TARS tackles a real, costly and repetitive problem. By turning translator assignment into a fast, consistent and explainable recommendation, it gives PMs back a large amount of time and removes much of the inconsistency of manual decisions, while grounding those decisions in the company's own history rather than in individual memory. For the company that means lower operational cost and more defensible decisions; for translators, work distributed on the basis of genuine fit and track record rather than on which manager thinks of them first, which is fairer. Throughout, the system assists rather than replaces: it offers a transparent starting point and leaves the final, accountable decision with a person.

\subsection{How We Could Have Made It Better}

We are clear-eyed about the system's limitations. We would ground the machine-learning model in real outcomes --- training it on the quality actually delivered or on whether tasks were on time, rather than on a designed efficiency score, so it discovers patterns the hand-set rules miss; our own analysis even points the way, since punctuality is the dimension that is actually predictable. We would evaluate at full scale, running the backtesting framework across the whole history to turn ``we think this works'' into evidence. We would simplify and unify the few places where the system grew two parallel ways of doing the same thing. We would add a fuller set of automated tests for the most intricate logic and harden the backend before any real deployment. And we would validate the business assumptions --- the resource-conservation choice and the exact relaxation and weighting numbers --- against what the company actually wants, ideally with PMs in the loop.

\subsection{What We Would Change About the Process}

If we started again, we would invest in the data-quality analysis even earlier and treat it explicitly as the foundation for everything else, since so many later decisions depended on understanding the data. We would decide sooner what precisely we wanted the model to predict, rather than building the modelling machinery first and questioning the target afterwards. And we would lock in the shared data formats and component boundaries earlier, to reduce the integration friction of keeping independently built parts in sync. None of these are surprising lessons --- they are the normal trade-offs of building something real under time pressure --- but they are the ones we felt most directly and would carry into the next project.

\subsection{Closing}

TARS demonstrates that a large, expensive, manual decision can be turned into a fast, consistent and explainable recommendation using a company's own historical data. The result is not a system that decides for people, but one that helps them decide better --- saving time, spreading work more fairly, and capturing institutional knowledge that would otherwise stay locked in individual heads. There is clear room to make it more accurate and more polished, but as a foundation it shows both the value of the idea and a realistic path to delivering it.

\section*{References}
\addcontentsline{toc}{section}{References}

\begin{enumerate}[label={[\arabic*]},leftmargin=*]
    \item The pandas development team. pandas: Python Data Analysis Library. \url{https://pandas.pydata.org}
    \item T. Chen and C. Guestrin. XGBoost: A Scalable Tree Boosting System. \url{https://xgboost.readthedocs.io}
    \item Scikit-learn developers. Scikit-learn: Machine Learning in Python. \url{https://scikit-learn.org}
    \item S. Ram\'irez. FastAPI. \url{https://fastapi.tiangolo.com}
    \item Meta Open Source. React. \url{https://react.dev}
\end{enumerate}

%----------------------------------------------------------

\end{document}
